\vfill\pagebreak%
\section{Operators}
\label{sec:types:operators}

A program that could only move values from box to box would not be worth much.
Operators are what build new values out of old ones, and an rvalue is usually not
a bare value but an expression: one or more variables combined by operators. They
fall into four groups.

\vspace{-5pt}%
\begin{itemize}
  \setlength\itemsep{-4pt}
\item \textbf{Arithmetic} operators -- addition, subtraction, multiplication and
  the rest -- apply to integers and floating-point numbers, and return a value of
  the same type as their operands.
\smallskip{}

\item \textbf{Comparison} operators relate two values and return a
  \token{bool}. They are what a program branches on.
\smallskip{}

\item \textbf{Bitwise} operators work on the individual bits of a value. They
  wait for a later chapter, because using them correctly means knowing first how
  a value is laid out in memory.
\smallskip{}

\item \textbf{Boolean} operators combine \token{bool} values according to
  Boolean logic.
  \smallskip{}

\end{itemize}

The listing below puts the first two groups to work on \token{i32} values.

\begin{lstlisting}[style=coloredverbatim, label=lst:types:operator_example, caption=Example of basic operators on \token{i32} values]
use std::io;

fn main() {
  let x = 12;
  let y = 24;

  let z = x + y;
  let w = z * 4;

  println("Z : ", z);
  println("W : ", w);
  let b = z > x;

  println("Z > X : ", b);
}
\end{lstlisting}
\vspace{-10pt}%
\begin{lstlisting}[style=bashVerb, label=lst:types:operator_example_output]
Z : 36
W : 144
Z > X : true
\end{lstlisting}

Several operators can appear in one expression, and then the question is which
runs first. As in mathematics, each operator has a precedence, and the higher
precedence goes first; Table~\ref{tab:types:operator_precedence} gives the order
for all of them.

Precedence is easiest to see as a tree. Each node is an operator, its children
are the operands it works on, and the leaves are the values written in the source
code. Evaluation runs from the leaves upwards, each node collapsing to a single
value, until the root is all that is left.

Figure~\ref{fig:types:operator_precedence} is the tree for
\token{1 + 3 * (4 - 9)}. Parentheses outrank every operator, so \token{4 - 9}
collapses first, even though \token{-} would otherwise yield to \token{*}.

\input{chapters/chapter2/figures/operator_precedence_table}

Figure~\ref{fig:types:operator_precedence_no_par} is the tree for
\token{1 + 3 * 4 - 9}, the same expression with the parentheses taken out. The
shape is different, and so is the answer: with nothing to override it, precedence
alone decides.

\input{chapters/chapter2/figures/operator_precedence_example}

The table also lists the compound assignments, \token{+=}, \token{-=} and their
kin. Each folds an operation into an assignment: the lvalue serves as the left
operand, the rvalue as the right, and the result goes back into the lvalue. So
\token{+=} adds to a variable and stores the sum in one step.

The rvalue is taken whole, whatever its precedence. In \token{x *= y + 1} the
addition binds more loosely than the multiplication, and it still happens first:
the statement means \token{x = x * (y + 1)}, not \token{x = x * y + 1}.

The compiler always knows which operator goes first. The reader of the code may
not, so in a long expression it is worth adding parentheses that change nothing,
or splitting the work over several lines.

\begin{lstlisting}[style=coloredverbatim]
let mut x = 10;
let y = 3;

x *= 2 + y; // rewritten into x = x * (2 + y);
println("X : ", x);
\end{lstlisting}
\vspace{-10pt}%
\begin{lstlisting}[style=bashVerb]
X : 50
\end{lstlisting}


Everything so far has been binary -- two operands. Unary operators take one, and
bind more tightly than any binary operator. Ymir has four: \token{-}, \token{&},
\token{*} and \token{!}. What each one means depends on the type it is applied
to:
\vspace{-8pt}%
\begin{itemize}
\setlength\itemsep{0.2mm}
\item \token{-} on an integer gives its opposite.
\item \token{-} does not apply to a \token{bool} at all.
\item \token{!} on a \token{bool} gives the other \token{bool} -- logical
  \textit{not}.
\item \token{!} on an integer is a bitwise operation, and waits for the chapter
  that covers those.
\end{itemize}

